\chapter{Electric Vehicles: The System-Level Win}
\label{ch:ev}

EVs are the capstone precedent because the automobile is not a component but a \emph{system}---the most complex mass-manufactured product in existence, wrapped in brands, dealer networks, safety regulation, and a century of incumbent advantage. If the playbook worked here, no consumer-industrial market is safe from it. It worked here.

\section{The leapfrog decision (2001--2012)}

The strategic insight predates Tesla. Wan Gang---an Audi engineer in Germany through the 1990s---persuaded Beijing around 2000 that China could never catch Toyota and VW in internal combustion, but could \emph{leapfrog} them if the industry changed energy carrier; EV powertrains entered the state 863 R\&D program in 2001, and Wan became Minister of Science and Technology in 2007~\cite{wan-gang}. The early demand instruments were clumsy and mostly failed---the 2009 ``Ten Cities, Thousand Vehicles'' pilots missed most targets; fewer than 500 EVs were sold nationally in 2009~\cite{ten-cities}. The state persisted anyway: NEVs were designated a Strategic Emerging Industry in 2010, and purchase subsidies, plate privileges, and procurement began compounding.

\section{The demand machine (2013--2022)}

Cumulative support to the sector, 2009--2023, is estimated by CSIS at \$230.9 billion---purchase subsidies, tax exemptions, infrastructure, and R\&D, a conservative figure excluding cheap credit and land~\cite{csis-230b}. The instruments were layered: purchase subsidies (wound down by end-2022, precisely as the market became self-sustaining); purchase-tax exemption extended through 2027; free ``green plates'' in cities where an ICE plate costs tens of thousands of dollars at auction; and the 2017 dual-credit mandate that forced every manufacturer, foreign JVs included, to produce NEVs or buy credits from those (BYD, Tesla) who did~\cite{dual-credit}. Beneath it all, infrastructure at a scale the West still does not comprehend: 19.3 million charging points by late 2025, of which 4.6 million public---roughly eighteen times the US public network~\cite{chargers}.

\begin{figure}[htbp]
\centering
\begin{tikzpicture}
\begin{axis}[bookaxis, height=5.4cm,
  xlabel={Year}, ylabel={Share of Chinese market (\%)},
  title={The fastest substitution event in automotive history},
  xmin=2019.6, xmax=2026.1, ymin=0, ymax=70,
  xtick={2020,2021,2022,2023,2024,2025},
  legend style={at={(0.02,0.98)}, anchor=north west},
]
\addplot[color=sOne, mark=*, mark size=1.6pt] coordinates {
  (2020,6.0) (2021,14.8) (2022,27.6) (2023,35.7) (2024,47.9) (2025,54.1)};
\addlegendentry{NEV retail penetration}
\addplot[color=sTwo, mark=square*, mark size=1.6pt] coordinates {
  (2020,62) (2021,56) (2022,50) (2023,44) (2024,39) (2025,35)};
\addlegendentry{foreign-brand share}
\node[font=\scriptsize\color{inkSec}, anchor=south east] at (axis cs:2025.9,55) {54.1\%};
\node[font=\scriptsize\color{inkSec}, anchor=north east] at (axis cs:2025.9,33) {35\%};
\draw[dashed, color=inkMut, line width=0.5pt] (axis cs:2023.4,0) -- (axis cs:2023.4,70);
\node[font=\scriptsize\color{inkSec}, anchor=north west, align=left]
  at (axis cs:2023.5,68) {crossover:\\NEVs pass\\foreign brands};
\end{axis}
\end{tikzpicture}
\caption[Chinese NEV penetration and foreign-brand collapse]{Two curves, one mechanism. New-energy-vehicle
retail penetration in China against the share of the Chinese market held by foreign brands, 2020--2025
[V]~\cite{cpca-2025,rhodium-evs}. December 2025 monthly NEV penetration passed 60\%. The technology transition
that reset incumbency and the collapse of incumbency are the same event, viewed from two sides.}
\label{fig:ev}
\end{figure}


The result is the fastest substitution event in automotive history (Figure~\ref{fig:ev}): NEV retail penetration went from $\sim$6\% (2020) to 27.6\% (2022) to 47.9\% (2024) to 54.1\% full-year 2025, crossing 60\% in December 2025~\cite{cpca-2025}. China has been the world's largest car market for seventeen consecutive years; in 2023 it passed Japan as the world's largest auto exporter, and exported 7.1 million vehicles in 2025 (2.6 million NEVs, doubling year-on-year); in June 2026 monthly exports passed one million for the first time, more than half NEVs~\cite{caam-2025,exports-june-2026}.

\section{BYD: the vertically integrated apex predator}

BYD is the pattern's exemplar. Battery maker first (Chapter~\ref{ch:batteries}), it acquired a moribund automaker in 2003 and spent fifteen years unremarkably; then the flywheel engaged: 416 thousand vehicles in 2020, 4.27 million in 2024, 4.55 million in 2025---taking the global BEV crown from Tesla by over 600,000 units~\cite{byd-2025}. The weapon is vertical integration: BYD makes approximately 80\% of its Tier-1 components in-house---cells, motors, power electronics including IGBT and SiC dies from BYD Semiconductor---avoiding an estimated \$2,369 per vehicle in supplier margins; Rhodium's teardown accounting puts BYD's structural cost advantage over Tesla at $\sim$\$4,700 per vehicle, of which direct subsidy is under \$300~\cite{rhodium-evs}. The Seagull retails under \$8,000 domestically. BYD even integrated logistics: an eight-ship ro-ro fleet, flagship capacity 9,200 vehicles, able to export a million cars a year~\cite{byd-ships}; and factories in Thailand, Brazil, Hungary, and Turkey place final assembly inside tariff walls.

Around the apex predator, an ecosystem of siege: Geely (3.0 million, owner of Volvo, Zeekr, and Lotus), Chery (1.34 million exported in 2025---one vehicle every 23 seconds), SAIC/MG (first Chinese brand past one million cumulative European sales), Xiaomi (289,000 orders for the YU7 in its first hour; 600,000 cumulative deliveries within 22 months of entering the car industry), XPeng, NIO, Li Auto, and Huawei's HIMA alliance~\cite{chery,mg-million,xiaomi}. Behind them, a graveyard: roughly 400 Chinese EV ventures died between 2018 and 2025~\cite{shakeout}. The state subsidizes the species, not the specimen.

\section{The incumbents' collapse and the reverse technology flow}

Foreign brands' share of the Chinese market fell from 62\% (2020) to 35\% (2025). GM, at 4.04 million Chinese sales in 2017, sold 1.8 million in 2024 and took over \$5 billion in China write-downs; VW's China deliveries fell from 4.2 million (2019) to 2.7 million (2025)~\cite{gm-collapse,vw-decline}. The humiliation is qualitative, not just quantitative---the technology now flows backward: VW paid \$700 million for 5\% of XPeng and will license XPeng's autonomous-driving stack for its China EVs from 2026; Audi builds on a SAIC platform; Stellantis paid €1.5 billion for 20\% of Leapmotor and distributes its cars; Toyota's China BEVs carry BYD powertrains~\cite{reverse-flow}. A century of licensing ran the other way; it took one product generation to reverse it.

Western trade response repeated the solar script, with the same result so far: US 100\% tariffs (May 2024) seal a market Chinese makers had barely entered; EU countervailing duties of 17--35\% (October 2024) are already being negotiated down toward a minimum-price regime while BYD's Hungarian, Turkish, and Thai plants tunnel under them~\cite{us-tariff,eu-cvd}.

\section{The EV as AI system: the bridge to Part II}

The EV war matters to this book's argument for one further reason: the automobile became a software and AI platform in Chinese hands. Chinese suppliers hold $\sim$95\% of global ADAS LiDAR shipments (Hesai, RoboSense); city-level navigate-on-autopilot deployed at scale in China from 2024; BYD fitted its ``God's Eye'' ADAS across its entire line down to the \$8,000 Seagull at zero additional cost~\cite{lidar,gods-eye}. McKinsey's 2026 assessment: Chinese vehicles are ``frequently more advanced than their Western counterparts'' in software and automated driving, developed on 20--24-month cycles against the West's 40--50, with $\sim$30\% lower BOM and capex~\cite{mckinsey-2026}. In other words: the first mass-market, safety-critical, embedded-AI product category is already Chinese-led---and it is the demand anchor for the domestic AI chips, LiDAR, and models of Part~II. Rare-earth motor magnets, where China's April 2025 export controls idled Ford and Suzuki lines within weeks, previewed the supply-chain leverage now being assembled in AI inputs~\cite{csis-rare-earth}.

\section{Lessons carried forward}

The EV war contributes the final lessons. \emph{Eighth}: pick the technology transition where incumbency resets to zero, and force the transition by regulation at home. \emph{Ninth}: a system-level product won end-to-end pulls its entire component stack to dominance with it (batteries, LiDAR, power semiconductors, magnets)---and the converse: dominance in components makes the next system-level product cheaper to win. \emph{Tenth}: speed is a weapon; a 2$\times$ development-cycle advantage compounds every generation. All ten lessons are now being applied to AI.
